\documentclass[11pt]{article}
\usepackage{preprint}

% authblk prints only the last \author when no \affil is given:
% list every author on one line instead (no affiliation marks).
\makeatletter
\renewcommand\AB@authnote[1]{}
\renewcommand\@author{\AB@authlist}
\makeatother
\DeclareRobustCommand{\authorstack}[1]{\begin{tabular}[t]{@{}c@{}}#1\end{tabular}}
\interfootnotelinepenalty=10000

\title{Addressing the Unique Academic Challenges of K--12 Students with Emotional and Behavioral Disorders}
\author{Zonghan Dai}
\author{Lavinia Wang}
\author{\authorstack{Jiadi He}}
\author{Harry Zhang}
\shorttitle{K--12 Students with Emotional and Behavioral Disorders}
\keywords{education policy; disability; implementation}
\runningauthors{Dai, Wang, He \& Zhang}
\papernote{Policy memo prepared for the U.S. Department of Education, Office of Special Education Programs.}

\begin{document}
\maketitle
\keywordsline

\section{Problem Statement}

How can the Office of Special Education Programs help K-12 students with EBD bridge the academic achievement gap with their peers?

Based on the criteria of cost, effectiveness and political feasibility, we recommend establishing a grant program to cultivate EBD specialist educators and ensure adequate staffing in every school.

\section{Defining the Problem}

In the United States, K-12 students with EBD (Emotional/Behavioral Disorders) face even more significant challenges in academic achievement after COVID-19, which causes them to lag far more behind their peers (Page, 2023). Their externalizing behaviors, such as aggression and destructiveness, along with internalizing behaviors like social withdrawal and anxiety, result in a lack of self-regulation and ability to learn effectively (Gidlund, 2018). COVID-19 has led to the transition of previously in-person interventions and support for this group to remote settings, which has further exacerbated the extent to which they fall behind.

Students with EBD classifications face unique challenges to learning in the mainstream classroom. It's important to emphasize that students with EBD are more susceptible to experiencing significant psychological issues, such as depression, anxious feelings, reduced self-confidence, and withdrawal from social interactions (Penn GSE Perspectives on Urban Education, n.d.). They may also face issues of low self-esteem and lack of self-awareness, they are deficient in social skills, unable to engage in sufficient validating interactions with peers and teachers. This can lead to their psychological needs not being met, affecting all aspects of their learning and life (Penn GSE Perspectives on Urban Education, n.d.). Also, teachers have responded that students with EBD are difficult to manage and communicate with because they cause some disciplinary disruptions to the teacher's instruction during up to 80\% of class time (Penn GSE Perspectives on Urban Education, n.d.)\footnote{Including students with emotional and behavioral disorders: Case management work protocol. Including Students with Emotional and Behavioral Disorders: Case Management Work Protocol | Penn GSE Perspectives on Urban Education. (n.d.). \url{https://urbanedjournal.gse.upenn.edu/archive/volume-20-issue-1-fall-2022/including-students-emotional-and-behavioral-disorders-case}}. Collectively, all these factors constitute the primary contributors to their academic underperformance (Mason \& Shriner, 2008).

\section{Alternatives}

\subsection{Status quo: Response to Intervention Program (RTI) and Individualized Education Program (IEP)}

The Federal Special Education Law of 2004 initially introduced the RTI program. The program's primary objective is to actively assist students with academic difficulties in improving their circumstances. Over time, the program has significantly enhanced its effectiveness and garnered widespread adoption in numerous schools. It has recognized students with diverse special educational needs as its primary beneficiaries (Education Week, 2020). The RTI program is exceptionally well-suited to actively support students experiencing emotional and behavioral challenges (EBD) or those at risk of developing them. Tier 1 revolves around delivering high-quality core instruction to proactively prevent risks. The overarching goals of Tier 1 prevention are more broadly focused, with the program's mission emphasizing the enhancement of the school's atmosphere, the promotion of collaboration between schools and families, and the provision of positive support to the target audience (Lehr, 2005).\footnote{Lehr, C. A. (2005). \textit{Impact: Volume 18, number 2: Students with emotional/behavioral disorders: Promoting positive outcomes}. Impact | Volume 18, Number 2 | Students with Emotional/Behavioral Disorders: Promoting Positive Outcomes | Institute on Community Integration Publications. \url{https://publications.ici.umn.edu/impact/18-2/students-with-emotional-behavioral-disorders}}

The Individualized Education Program (IEP) process represents a significant step in delivering services and enhancing the results for students with disabilities (Massachusetts Department of Education, 2001).\footnote{IEP Process Guide - Massachusetts Department of Elementary and Education (2001). \url{https://www.doe.mass.edu/sped/iep/proguide.pdf}} After evaluating the student's current academic level and identifying their learning requirements, the IEP team can proceed to create tailored instructional methods regarding special education, along with the necessary support (GreatSchools, 2023).\footnote{Aligning the IEP and academic content standards to improve academic achievement - GreatSchools (2023). \url{https://www.greatschools.org/gk/articles/iep-academic-content-standards/}}

\subsection{Establish a Grant Program to Cultivate EBD Specialist Educators: Ensure Adequate Staffing in Every School}

In the pursuit of inclusive and equitable education, nationwide schools must implement comprehensive policies addressing the unique needs of students with Emotional/Behavioral Disorders (EBD). This entails equipping each school with targeted programs, led by qualified EBD educators, to enhance academic outcomes. These experts, with a deep understanding of EBD intricacies, employ effective strategies for behavior management, emotion regulation, academic support, and overall growth. The policy advocates for a minimum of 1-2 qualified EBD educators and support staff in every school, ensuring individualized care and a psychologically safe environment for both students and parents.

Nationally, data from the School Mental Health Report Card in 2020 revealed that nearly 31 percent of parents believe their child's emotional well-being has deteriorated since the onset of the pandemic. Furthermore, in 2020, there was a 24\% increase in emergency room visits for mental health-related emergencies among children aged 5-11 and a 31\% increase among those aged 12-17. This challenging situation has persisted post-pandemic (Snapshot, 2022).\footnote{Snapshot: Student Mental Health Incidence and availability of therapeutic staff. NJCommonGround.org. (2022). \url{https://njcommonground.org/snapshot-student-mental-health-incidence-and-availability-of-therapeutic-staff/}} A 2021 report by the Substance Abuse and Mental Health Services Administration (SAMHSA) indicated that 14\% of 13-17 year olds experience mood disorders, and 19\% grapple with behavioral disorders. These statistics underscore the growing demand for specialized assistance related to EBD among an increasing number of K-12 students (New Jersey Department of Education, n.d.).\footnote{New Jersey Department of Education. Mental Health. (n.d.). \url{https://www.nj.gov/education/safety/wellness/mh/}}

\subsection{Provide grants to support the nationwide integration of SEL into K-12 curriculum design}

Social Emotional Learning programs have the potential to positively impact the development and academic performance of students with EBD. A study conducted by Espelage et al. in 2016 investigated the effects of attending a school with an SEL program for three years on students with disabilities, focusing on their prosocial behaviors and academic achievements (Hamline University, 2021). This policy alternative strongly advocates the widespread integration of SEL programs throughout K-12 classrooms. This integration entails the implementation of various activities, such as collaborative teamwork exercises, role-playing involving both teachers and students, and leadership training. For instance, it recommends infusing SEL-related lessons into traditional subjects like math, history, or reading. In practical terms, this approach might involve students voluntarily assuming roles and collaborating on group projects, immersing themselves in role-play scenarios of historical figures to gain insight into human behavior, or conducting formal interviews with their peers to stay informed about current events (National University, 2022).\footnote{National University. (2022). \textit{What is social emotional learning (SEL): Why it matters}. \url{https://www.nu.edu/blog/social-emotional-learning-sel-why-it-matters-for-educators/}}

\section{Criteria}

\noindent Cost:
\begin{itemize}
\item High: Costs are deemed high if the alternative necessitates substantial financial investment, such as the development of new infrastructure or a significant shift in educational curriculum requiring extensive resources.
\item Medium: Costs are categorized as moderate when the alternative predominantly utilizes existing structures, staff, and resources, minimizing the need for major financial commitments.
\item Low: The policy is considered to have low costs when it requires minimal levels of investment, utilizing existing resources efficiently without significant financial implications.
\end{itemize}

\noindent Effectiveness:
\begin{itemize}
\item High: The implementation of policy alternatives have been extended to EBD students nationwide, and these students' test scores have improved significantly.
\item Medium: Policy alternatives are trying to reach all students with EBD across the country, but they are not mandatory in all states or school districts, and not all test scores for students with EBD have improved, nor has there been particularly much positive feedback from parents.
\item Low: The districts in which the policy alternatives were implemented are a small percentage of the nation, and there are insufficient quantitative data results to indicate that the gap between EBD students' test scores and the rest of their peers has narrowed.
\end{itemize}

\noindent Feasibility:
\begin{itemize}
\item High: Feasibility is rated as high when the alternative is easily implementable, requiring minimal logistical challenges, and the necessary resources are readily available, allowing for a smooth nationwide adoption.
\item Medium: Feasibility is considered moderate when the alternative faces some implementation challenges, such as varied regulations across states or districts. The required resources may be available, but coordination efforts are needed for a more widespread application.
\item Low: Feasibility is deemed low when the alternative encounters significant implementation obstacles, requiring substantial coordination, adaptation, and potential resistance. Limited availability of essential resources may impede widespread adoption.
\end{itemize}

\section{Discussion of Criteria}

The criteria for evaluating the proposed alternatives are key in guaranteeing they align with the broader goals and values that motivate our work in addressing the academic achievement gap for K-12 students with EBD. Our objective to promote inclusion, fairness, and holistic rehabilitating for K-12 students struggling with emotional/behavioral disorders (EBD) is based on three key principles: cost, effectiveness, and feasibility. These standards embody our unshakable commitment to guaranteeing that every student with EBD has an equal opportunity for academic success and all-encompassing assistance.

The cost factor is the basis that directly affects how resources are allocated, which is essential for the implementation of a strong grant program that develops EBD expert educators. It becomes essential to strike a balance between financial availability and resource optimization to successfully educate and position these professionals in all educational settings.\footnote{Response to intervention 2.0. Education Week. (2020, November 10). \url{https://www.edweek.org/teaching-learning/response-to-intervention-2-0}}

Effectiveness becomes the primary focus in this framework, essential to its goal of providing individualized and customized support to kids struggling with EBD. The program's effectiveness will be evaluated based on how effectively it affects these children's academic progress, emotional health, and overall outcomes progression.

A key role is played by feasibility, which captures the practical viability of a grant program countrywide. This criterion includes the complex web of logistical difficulties, the availability of resources, and the initiative's flexibility in a range of educational contexts. Securing the triumph of this endeavor in many settings continues to be the central goal of our wider goals and principles.

\section{Outcomes}

\subsection{Alternative 1: Status Quo (RTI and IEP)}

\begin{itemize}
\item Cost: The status quo alternatives, RTI and IEP, often rely on existing structures and resources, resulting in a lower financial investment compared to implementing new programs.\footnote{Response to intervention 2.0. Education Week. (2020, November 10). \url{https://www.edweek.org/teaching-learning/response-to-intervention-2-0}} However, ongoing resource allocation for training and specialized support might be necessary to maintain effectiveness.
\item Effectiveness: Both RTI and IEP have demonstrated positive outcomes in providing tailored interventions for students with EBD, albeit with varying degrees of success across different regions and schools. They contribute to improved academic achievements and individualized assistance for students.\footnote{Lehr, C. A. (2005). Impact: Volume 18, number 2: Students with emotional/behavioral disorders: Promoting positive outcomes. Impact | Volume 18, Number 2 | Students with Emotional/Behavioral Disorders: Promoting Positive Outcomes | Institute on Community Integration Publications. \url{https://publications.ici.umn.edu/impact/18-2/students-with-emotional-behavioral-disorders}}
\item Feasibility: These established programs, while effective, may pose fewer feasibility challenges due to their familiarity with educational systems. However, ensuring consistent implementation and adherence to standards across diverse regions could present moderate challenges.
\end{itemize}

\subsection{Alternative 2: Establish a Grant Program to Cultivate EBD Specialist Educators}

\begin{itemize}
\item Cost: The establishment of a grant program to cultivate EBD specialist educators may require a moderate to high initial financial investment. This investment will cover training, staffing, and ongoing support but is expected to yield significant long-term benefits for students with EBD.\footnote{Page, Dana, ``A meta-analysis of the academic achievement of students with emotional/behavioral disorders in traditional public schools in the United States.'' (2023). Electronic Theses and Dissertations. Paper 4085. Retrieved from \url{https://ir.library.louisville.edu/etd/4085}}
\item Effectiveness: This alternative is projected to have a positive impact on the effectiveness criterion, as specialized EBD educators can provide tailored support leading to improved academic outcomes and emotional well-being for students.\footnote{Natasha Barnett. (2023). Effective assessment for exceptional Ebd students. Share \& Discover Presentations | SlideShare. \url{https://www.slideshare.net/natashabarnettmiamig/effective-assessment-for-exceptional-ebd-students}}
\item Feasibility: Implementing a nationwide grant program to cultivate EBD specialist educators might pose moderate feasibility challenges due to logistical considerations, such as staffing every school with qualified specialists and coordinating resources on a national scale.
\end{itemize}

\subsection{Alternative 3: Provide Grant to Implement Nationwide SEL Integration}

\begin{itemize}
\item Cost: Integrating Social Emotional Learning (SEL) into the curriculum might involve moderate to high initial costs for training educators and developing specialized materials. However, the long-term benefits could outweigh these costs by improving students' emotional well-being and academic performance.
\item Effectiveness: SEL programs have shown promise in fostering positive social and emotional skills in students, including those with EBD. However, the effectiveness may vary based on the comprehensive integration of SEL into diverse educational settings.
\item Feasibility: Implementing SEL programs nationwide may pose moderate to high feasibility challenges due to the need for extensive coordination among educators, development of specialized curriculum, and varying state or district regulations.
\end{itemize}

\section{Analysis Matrix}

\begin{center}
\renewcommand{\arraystretch}{1.15}
\begin{tabularx}{\textwidth}{@{}>{\raggedright\arraybackslash}p{0.17\textwidth}*{3}{>{\raggedright\arraybackslash}X}@{}}
\toprule
\textbf{Criteria} & \textbf{Alternative 1: Status Quo (RTI and IEP)} & \textbf{Alternative 2: Establish a Grant Program to Cultivate EBD Specialist Educators} & \textbf{Alternative 3: Provide Grant to Implement Nationwide SEL Integration} \\
\midrule
\textbf{Cost} & High & Medium & High \\
\textbf{Effectiveness} & Low & High & Medium \\
\textbf{Feasibility} & High & High & Low \\
\bottomrule
\end{tabularx}
\end{center}

Choosing Alternative 2, Cultivate EBD Specialist Educators, emerges as the best option for supporting students with Emotional and Behavioral Disorders (EBD). This approach, despite requiring a higher initial investment, promises greater long-term benefits. It excels in effectiveness, as specialized educators can provide more personalized and impactful support for students' academic and emotional needs. While there are challenges in implementing this nationwide, such as ensuring qualified staffing and resource coordination, these are manageable with proper planning. This alternative stands out for its direct, focused approach to improving outcomes for students with EBD, compared to the other options.

\section{Projection}

The projections for cost, effectiveness, and feasibility stemmed from a meticulous analysis of the proposed policies. For instance, cost projections were grounded in the anticipated requirements of establishing a grant program to foster EBD specialist educators. This encompassed factors like initial training expenses, ongoing support, and the staffing needs across schools. Consideration was given to strike a balance between investing adequately for impactful outcomes and ensuring cost efficiency.

Effectiveness projections relied on the anticipated positive influence of specialized support offered by EBD educators on students' academic progress and emotional well-being. This outlook was drawn from existing evidence showcasing the potential benefits of tailored assistance for students with EBD, emphasizing improved academic performance and better emotional regulation.

Feasibility projections were derived from acknowledging the logistical intricacies inherent in implementing a nationwide program. Recognizing the challenges of staffing every school with qualified specialists and coordinating resources on a national scale underscored the feasibility concerns. Balancing the aspirations of comprehensive support with the pragmatic understanding of logistical hurdles formed the foundation of these feasibility projections.

The integration of these projections considered the intricate interplay between financial, effectiveness, and logistical factors within the proposed framework, aiming to strike a balance between ambition and practical implementation realities.

\section{Tradeoffs}

While the status quo, represented by the Response to Intervention Program (RTI) and Individualized Education Program (IEP), has proven effective in addressing academic difficulties and providing tailored services for students with disabilities, it may fall short in meeting the specific and growing needs of students with Emotional/Behavioral Disorders (EBD). RTI focuses on prevention through high-quality core instruction but may not provide sufficient specialized support for EBD students, potentially leading to unmet mental health challenges.

In contrast, the proposed Grant Program to Cultivate EBD Specialist Educators offers a targeted approach, acknowledging the rising demand for specialized assistance in emotional and behavioral areas. This recommendation emphasizes the importance of qualified EBD educators and support staff in every school, promising individualized care and a psychologically safe environment. The tradeoff involves redirecting resources from broad programs like RTI and IEP to a more specialized initiative, ensuring a comprehensive response to the growing mental health challenges faced by students, particularly those with EBD.

\section{Uncertainty}

Budget availability and allocation are crucial factors in the efficacy of the three alternatives discussed in the analysis for aiding students with Emotional and Behavioral Disorders (EBDs). The feasibility of implementing Response to Intervention (RTI) and Individualized Education Programs (IEPs), cultivating EBD specialist educators, or integrating Social-Emotional Learning (SEL) into curricula hinges on having adequate financial resources. Insufficient funding could significantly impede these initiatives, leading to a shortage of trained personnel and inadequate monitoring of outcomes.

Beyond financial considerations, the success of these alternatives also depends on support from parents and the community. Active parental involvement and community engagement can enhance the effectiveness of interventions for children with EBD, extending support beyond the school environment. Informed and engaged parents can bolster school-based initiatives and provide essential feedback. Moreover, community attitudes towards EBD can influence stigma and the level of inclusion within different areas. A supportive and collaborative atmosphere from families and communities is crucial for both students and educators. Nonetheless, there is variability in the level of support across school districts, which adds another layer of uncertainty to the analysis. To simplify the evaluation, it's assumed that parental and community support is uniformly positive, though this may not reflect the complex realities in different educational settings.

Adjusting the assumptions about budget availability and parental/community support could significantly alter the outcomes of the three alternatives, evaluated against effectiveness, cost, and political feasibility.

First Alternative (Status Quo - RTI and IEP): With limited budgets, the effectiveness of RTI and IEP programs could diminish. Restricted financial resources might lead to reduced training and resources, essential for effective implementation. However, the cost in terms of financial investment would not change significantly, nor would the political feasibility, as these are established programs with existing support.

Second Alternative (EBD Specialist Educators): A budget cut would likely have a moderate impact on the effectiveness of cultivating EBD specialist educators. Limited funding could hinder the depth and quality of training and support, crucial for the program's success. While the direct cost of the program might not be significantly affected, reduced funding could influence stakeholder support, including from policymakers, potentially impacting political feasibility.

Third Alternative (Nationwide SEL Integration): Insufficient budgeting would heavily impact the effectiveness of integrating SEL into curricula. The development and implementation of a new curriculum requires substantial investment. While the cost in terms of budget allocation remains high, political feasibility might not change significantly, but the program's effectiveness would be compromised due to financial constraints.

Impact of Parental and Community Involvement: Reduced parental and community involvement would affect all three alternatives. Lower engagement could lead to less effective outcomes, increasing the long-term cost and potentially affecting political feasibility. This is especially relevant for students with EBD who benefit significantly from sustained support and involvement from their community and family networks. The effectiveness of all programs is somewhat contingent on this external support, highlighting the importance of community and parental engagement in educational initiatives.

\section{Recommendation}

Based on the practical considerations and the real-world impact, we recommend establishing a grant program to cultivate EBD specialist educators and ensuring adequate staffing in every school. It not only addresses the academic achievement gap for students with EBD but also aligns with the values of inclusivity, equality, and positive outcomes. The investment in specialized educators reflects a commitment to the well-being and success of students with EBD, contributing to a more inclusive and supportive educational system.

\begin{references}[Reference]
\item Barni, D., Danioni, F., \& Benevene, P. (2019). Teachers’ self-efficacy: The role of personal values and motivations for teaching. \textit{Frontiers in psychology}, \textit{10}, 1645.
\item Carmo Cabral-Gouveia. (2023). \textit{Educational strategies to reduce the achievement gap: a systematic review}. Bing. \url{https://www.frontiersin.org/articles/10.3389/feduc.2023.1155741/full}
\item Cooc, N. (2017). Examining Racial Disparities in Teacher Perceptions of Student Disabilities. Teachers College Record, 119(7), 1-32. \url{https://doi.org/10.1177/016146811711900703}
\item Effective programs for emotional and behavioral disorders - HSD. (2013). \url{https://psych.hsd.ca/Effective%20Programs%20for%20EBDs.pdf}
\item Emotional/Behavioral Disorders: Promoting Positive Outcomes | Institute on Community Integration Publications. \url{https://publications.ici.umn.edu/impact/18-2/students-with-emotional-behavioral-disorders}
\item Gamage, K. A. A., Dehideniya, D. M. S. C. P. K., \& Ekanayake, S. Y. (2021). The Role of Personal Values in Learning Approaches and Student Achievements. \textit{Behavioral sciences (Basel, Switzerland)}, \textit{11}(7), 102. \url{https://doi.org/10.3390/bs11070102}
\item Gidlund, U. (2018). Teachers’ attitudes towards including students with emotional and behavioural difficulties in mainstream school: A systematic research synthesis. International Journal of Learning, Teaching and Educational Research, 17(2), 45-63. \url{https://doi.org/10.26803/ijlter.17.2.3}
\item Hamline University. (2021). What Social-emotional Learning Looks Like In The High School Special Education Setting, And How It Benefits Students With Emotional Or Behavioral Disorders. \url{https://digitalcommons.hamline.edu/cgi/viewcontent.cgi?article=1658&context=hse_cp}
\item \textit{Including students with emotional and behavioral disorders: Case management work protocol}. Including Students with Emotional and Behavioral Disorders: Case Management Work Protocol | Penn GSE Perspectives on Urban Education. (n.d.). \url{https://urbanedjournal.gse.upenn.edu/archive/volume-20-issue-1-fall-2022/including-students-emotional-and-behavioral-disorders-case}
\item John McKenna Associate Professor of Special Education. (2023,). Students with disabilities are not getting help to address lost opportunities. The Conversation. \url{https://theconversation.com/students-with-disabilities-are-not-getting-help-to-address-lost-opportunities-177431}
\item Kivimäki V, Meriluoto S. Holistic Perspective to Individual Study Plan: Personal Development Project Plan. Stud Health Technol Inform. 2018;256:514-526. PMID: 30371411.
\item Lehr, C. A. (2005). \textit{Impact: Volume 18, number 2: Students with emotional/behavioral disorders: Promoting positive outcomes}. Impact | Volume 18, Number 2 | Students with Emotional/Behavioral Disorders: Promoting Positive Outcomes | Institute on Community Integration Publications. \url{https://publications.ici.umn.edu/impact/18-2/students-with-emotional-behavioral-disorders}
\item Lovell A, Bailey J. Nurses' perceptions of personal attributes required when working with people with a learning disability and an offending background: a qualitative study. J Psychiatr Ment Health Nurs. 2017 Feb;24(1):4-14. doi: 10.1111/jpm.12326. Epub 2016 Jul 21. PMID: 27439391.
\item Mason, \& Shriner, J. G. (2008). Self-regulated strategy development instruction for writing an opinion essay: Effects for six students with emotional/behavior disorders. Reading \& Writing, 21(1-2), 71--93. \url{https://doi.org/10.1007/s11145-007-9065-y}
\item Natasha Barnett. (2023). \textit{Effective assessment for exceptional Ebd students}. Share \& Discover Presentations | SlideShare. \url{https://www.slideshare.net/natashabarnettmiamig/effective-assessment-for-exceptional-ebd-students}
\item \textit{New Jersey Department of Education}. Mental Health. (n.d.). \url{https://www.nj.gov/education/safety/wellness/mh/}
\item Okonofua, J. A., \& Eberhardt, J. L. (2015). Two Strikes: Race and the Disciplining of Young Students. Psychological Science, 26(5), 617-624. \url{https://doi.org/10.1177/0956797615570365}
\item Page, Dana, ``A meta-analysis of the academic achievement of students with emotional/behavioral disorders in traditional public schools in the United States.'' (2023). Electronic Theses and Dissertations. Paper 4085. Retrieved from \url{https://ir.library.louisville.edu/etd/4085}
\item Reid, Gonzalez, J. E., Nordness, P. D., Trout, A., \& Epstein, M. H. (2004). A Meta-Analysis of the Academic Status of Students with Emotional/Behavioral Disturbance. The Journal of Special Education, 38(3), 130--143. \url{https://doi.org/10.1177/00224669040380030101}
\item \textit{Response to intervention 2.0}. Education Week. (2020, November 10). \url{https://www.edweek.org/teaching-learning/response-to-intervention-2-0}
\item \textit{Snapshot: Student Mental Health Incidence and availability of therapeutic staff}. NJCommonGround.org. (2022). \url{https://njcommonground.org/snapshot-student-mental-health-incidence-and-availability-of-therapeutic-staff/}
\item Stanford University. (2016). \textit{Teachers more likely to label black students as Troublemakers, stanford research shows}. Stanford News. \url{https://news.stanford.edu/2015/04/15/discipline-black-students-041515/}
\item U.S. Department of Education. (n.d.). \textit{Implementation and Effectiveness of the Response to Intervention (RTI) Program}. ERIC - Education Resources Information Center. \url{https://files.eric.ed.gov/fulltext/EJ1099646.pdf}
\item Wynne, M. E., Ausikaitis, A. E., \& Satchwell, M. (2013). Adult outcomes for children and adolescents with EBD. \textit{SAGE Open}, \textit{3}(1), 215824401348313. \url{https://doi.org/10.1177/2158244013483133}
\end{references}

\end{document}
